\subsection{Adaptive borrowing with estimated verification probabilities}
\label{sec:supp-pi-adaptive}

Proposition~\ref{prop:pi-extension} treats propensity replacement during process evaluation for a fixed cross-fitted augmentation. The adaptive procedure also uses the verification probabilities in the gain criterion and therefore requires control of the selected coefficient.

For evaluation fold \(k\), enlarge \(\mathcal T_k\) to include all data used to fit \(\widehat q_k\), \(\widehat m_k\), and \(\widehat\pi_k\). The propensity regression is fitted outside the gain and evaluation folds. For a verification function \(r\), let
\[
(\widehat{\mathcal G}_{h,k}(r),
 \widehat{\mathcal H}_{h,k}(r),
 \widehat\lambda_{h,k}(r))
\]
denote the gain quantities computed with \(r\) in both the transfer weight and the inverse-probability residual. We compare \(r=\pi_M\) and \(r=\widehat\pi_k\), holding the threshold regressions and selection-fold observations fixed. Write
\[
\widehat a_{k,y}(\lambda)
=
\widehat q_{k,y}
+
\lambda\widehat D_{k,y},
\qquad
\widehat D_k=\widehat m_k-\widehat q_k,
\]
and set
\[
\widehat a_k[r]
=
\widehat a_k\{\widehat\lambda_{h,k}(r)\}.
\]

\begin{proposition}[Adaptive process with estimated verification probabilities]
\label{prop:pi-adaptive}
Suppose the conditions of Theorems~\ref{thm:process} and \ref{thm:bootstrap} hold for the known-design procedure. Assume, uniformly over the three rotations,
\begin{equation}
|\widehat{\mathcal G}_{h,k}(\widehat\pi_k)
 -\widehat{\mathcal G}_{h,k}(\pi_M)|
+
|\widehat{\mathcal H}_{h,k}(\widehat\pi_k)
 -\widehat{\mathcal H}_{h,k}(\pi_M)|
=o_{\mathbb P}(1).
\label{eq:pi-gain-perturbation}
\end{equation}
Assume Conditions~\eqref{eq:pi-drift-cond}--\eqref{eq:pi-var-cond} hold uniformly over
\[
\{\widehat a_k(\lambda):0\le\lambda\le1\}.
\]
For the coefficient direction, assume
\begin{align}
&\sqrt{N_{M,h}}
\sup_{y\in\mathcal Y}
\left|
\frac1{\mu_h}
E\left[
K_h\widehat D_{k,y}(W)
\frac{\widehat\pi_k(W)-\pi_M(W)}{\widehat\pi_k(W)}
\Bigm|\mathcal T_k
\right]
\right|
=o_{\mathbb P}(1),
\label{eq:pi-D-drift}\\
&\sup_{y\in\mathcal Y}
\frac{p_Mh^d}{\mu_h^2}
E\left[
K_h^2R\widehat D_{k,y}^2(W)
\left\{\frac1{\widehat\pi_k(W)}-\frac1{\pi_M(W)}\right\}^2
\Bigm|\mathcal T_k
\right]
=o_{\mathbb P}(1).
\label{eq:pi-D-variance}
\end{align}
Assume the propensity-induced evaluation, gain, and direction classes satisfy Condition~\ref{cond:entropy}. Then
\begin{equation}
\sqrt{N_{M,h}}
\left\|
\widehat F_h[\widehat\pi]
-
\widehat F_h[\pi_M]
\right\|_{\mathcal Y}
=o_{\mathbb P}(1).
\label{eq:pi-adaptive-equivalence}
\end{equation}
The Gaussian-process and multiplier limits coincide with those of the known-design estimator.
\end{proposition}

\begin{proof}
For fold \(k\), decompose
\begin{align}
\widehat F_{h,k}[\widehat\pi_k]-\widehat F_{h,k}[\pi_M]
&=
\{\widehat F_{h,k}(\widehat a_k[\widehat\pi_k],\widehat\pi_k)
  -\widehat F_{h,k}(\widehat a_k[\pi_M],\widehat\pi_k)\}
\label{eq:pi-coefficient-component}\\
&\quad+
\{\widehat F_{h,k}(\widehat a_k[\pi_M],\widehat\pi_k)
  -\widehat F_{h,k}(\widehat a_k[\pi_M],\pi_M)\}.
\label{eq:pi-evaluation-component}
\end{align}
For \eqref{eq:pi-evaluation-component}, Lemma~\ref{lem:propensity-drift}, Conditions~\eqref{eq:pi-drift-cond}--\eqref{eq:pi-var-cond}, and the conditional maximal inequality give a uniform \(o_{\mathbb P}(N_{M,h}^{-1/2})\) bound.

For \eqref{eq:pi-coefficient-component},
\begin{equation}
U_y(\widehat a_k[\widehat\pi_k],\widehat\pi_k)
-
U_y(\widehat a_k[\pi_M],\widehat\pi_k)
=
d_{\lambda,k}
\widehat D_{k,y}(W)
\left(1-\frac R{\widehat\pi_k(W)}\right),
\label{eq:pi-coefficient-signal}
\end{equation}
where
\[
d_{\lambda,k}
=
\widehat\lambda_{h,k}(\widehat\pi_k)
-
\widehat\lambda_{h,k}(\pi_M).
\]
Its localized conditional mean is
\begin{equation}
\frac1{\mu_h}
E\left[
K_h\widehat D_{k,y}(W)
\left\{1-\frac{\pi_M(W)}{\widehat\pi_k(W)}\right\}
\Bigm|\mathcal T_k
\right].
\label{eq:pi-coefficient-drift}
\end{equation}
Condition~\eqref{eq:pi-D-drift} makes the scaled mean contribution negligible because \(|d_{\lambda,k}|\le1\).

For the centered component, use
\[
1-\frac R{\widehat\pi_k}
=
\left(1-\frac R{\pi_M}\right)
+
R\left(\frac1{\pi_M}-\frac1{\widehat\pi_k}\right).
\]
The first term has localized squared radius
\[
\sup_y
\frac{p_Mh^d}{\mu_h^2}
E\left[
K_h^2
\left(\frac1{\pi_M}-1\right)
\widehat D_{k,y}^2
\Bigm|\mathcal T_k
\right],
\]
and Condition~\eqref{eq:pi-D-variance} makes the centered second term negligible. If \(\mathcal H_{h,k}\to H_{\infty,k}>0\), consistency of the known-design gain quantities, \eqref{eq:pi-gain-perturbation}, and local Lipschitz continuity of the projected ratio give
\[
d_{\lambda,k}=o_{\mathbb P}(1),
\]
while the first radius is \(O_{\mathbb P}(1)\). If \(\mathcal H_{h,k}\to0\), Condition~\ref{cond:path-compatibility} makes this radius tend to zero, while the coefficient difference remains bounded by one. The conditional maximal inequality and the ratio linearization therefore give
\[
\sqrt{N_{M,h}}
\left\|
\widehat F_{h,k}(\widehat a_k[\widehat\pi_k],\widehat\pi_k)
-
\widehat F_{h,k}(\widehat a_k[\pi_M],\widehat\pi_k)
\right\|_{\mathcal Y}
=o_{\mathbb P}(1).
\]
Summing over the evaluation folds proves \eqref{eq:pi-adaptive-equivalence}.

For the multiplier process, the estimated-propensity minus known-design influence difference has the same mean-plus-centered decomposition. The drift conditions remove its empirical centering at the \(N_{M,h}^{-1/2}\) scale, and the two variance conditions give a uniform empirical \(L_2\) norm of order \(o_{\mathbb P}(1)\). Lemma~\ref{lem:IF-replacement} and the conditional multiplier maximal inequality then transfer the known-design multiplier limit.
\end{proof}

A lower-level route uses local uniform relative consistency. Let \(C_K\) satisfy \(K(u)=0\) for \(\|u\|>C_K\), and define
\[
\epsilon_{\pi,M,h}
=
\operatorname*{ess\,sup}_{\|X-x_0\|\le C_Kh}
\left|\frac{\widehat\pi_k(W)}{\pi_M(W)}-1\right|.
\]
If \(\epsilon_{\pi,M,h}=o_{\mathbb P}(1)\), \(\widehat\pi_k/\pi_M\) is bounded above and away from zero on the local region, and the candidate paths are bounded, then direct expansion of the normalized gain integrands gives
\begin{equation}
|\widehat{\mathcal G}_{h,k}(\widehat\pi_k)
 -\widehat{\mathcal G}_{h,k}(\pi_M)|
+
|\widehat{\mathcal H}_{h,k}(\widehat\pi_k)
 -\widehat{\mathcal H}_{h,k}(\pi_M)|
=O_{\mathbb P}(\epsilon_{\pi,M,h}).
\label{eq:pi-primitive-gain}
\end{equation}
Define
\[
r_{D,M,h}^2
=
\sup_y
\frac1{\mu_h}
E\{K_h\widehat D_{k,y}^2\mid\mathcal T_k\}.
\]
The same relative-consistency condition makes the left side of \eqref{eq:pi-D-variance} \(O_{\mathbb P}(\epsilon_{\pi,M,h}^2)\). Moreover,
\[
\sqrt{N_{M,h}}\,r_{D,M,h}r_{\pi,M,h}
=o_{\mathbb P}(1)
\]
controls \eqref{eq:pi-D-drift}, while
\[
\sqrt{N_{M,h}}\,r_{a,M,h}r_{\pi,M,h}
=o_{\mathbb P}(1)
\]
controls the evaluation drift uniformly over the fitted segment. Together with \eqref{eq:pi-primitive-gain}, these conditions yield the adaptive process and multiplier conclusions.
